\ifdefined\gkver\else\def\gkver{student}\fi
\documentclass[\gkver]{gaokaozhenti}
\begin{document}

\chapter{2026年广东}

\begin{enumerate}
\item
%% number: 1
%% paperName: 2026年广东高考第 1 题 4 分
%% typeId: 1
%% type: 单选题
%% chapter: 原子和原子核
%% point: 玻尔模型
%% method: 无
%% score: 4
%% degree: 850
%% duplicateId: 0
%% body:
核钟是基于原子核能级跃迁来建立高精度时间标准的装置，原子核可吸收光子发生类似原子的能级跃迁。现有一种激光能够激发某原子核从基态跃迁至激发态，其能级差约为 $1.3\times10^{-18}\UJ$，普朗克常量 $h=6.6\times10^{-34}\UJcdots$，真空中光速 $c=3.0\times10^{8}\Ums$。关于该激光，下列说法正确的是 \xzanswer{D}

\fourchoices[answer=D]
{光强倍增，此能级跃迁不能发生}
{光强减半，此能级跃迁不能发生}
{频率约为 $2.0\times10^{16}\UHz$}
{波长约为 $1.5\times10^{-7}\Um$}

%% answer: D
%% memo:
\memoanswer{
A．原子核能级跃迁的条件是入射光子能量等于能级差，与光强无关，光强倍增只要光子能量匹配仍可发生跃迁，故 A 错误；

B．同理，光强减半不影响光子能量，只要能量匹配仍能发生跃迁，故 B 错误；

C．根据光子能量公式 $\Delta E=h\nu$，可得频率 $\nu=\frac{\Delta E}{h}=\frac{1.3\times10^{-18}}{6.6\times10^{-34}}\UHz\approx2.0\times10^{15}\UHz$，与选项给出的 $2.0\times10^{16}\UHz$ 不符，故 C 错误；

D．根据波长与频率关系 $\lambda=\frac{c}{\nu}=\frac{hc}{\Delta E}$，代入数值计算得 $\lambda=\frac{6.6\times10^{-34}\times3.0\times10^{8}}{1.3\times10^{-18}}\Um\approx1.5\times10^{-7}\Um$，故 D 正确。

故选 D。
}

\item
%% number: 2
%% paperName: 2026年广东高考第 2 题 4 分
%% typeId: 1
%% type: 单选题
%% chapter: 电场
%% point: 电势能电势差电场力做功
%% method: 无
%% score: 4
%% degree: 750
%% duplicateId: 0
%% body:
静电卡盘是芯片制造中的重要设备，如图为双极型静电卡盘吸附原理简图，双电极接高压电源后，晶片靠近卡盘的一侧带上与电极极性相反的电荷，在电场力作用下向卡盘运动并被吸附。在晶片运动过程中，下列说法正确的是 \xzanswer{B}
\onepicture[w=4.5cm]{figs/02.png}

\fourchoices[answer=B]
{电场力对晶片做负功}
{晶片在电场中的电势能逐渐减少}
{晶片与卡盘电极之间表现为斥力}
{晶片与卡盘电极之间电场强度处处相同}

%% answer: B
%% memo:
\memoanswer{
AC．晶片在电场力作用下向卡盘运动并被吸附，则晶片与卡盘电极之间表现为引力，晶片所受电场力方向指向卡盘，与晶片向卡盘运动的位移方向同向，电场力对晶片做正功，AC 错误；

B．电场力对晶片做正功，则晶片在电场中的电势能逐渐减少，故 B 正确；

D．根据题意可知该电场不是匀强电场，晶片与卡盘电极之间电场强度并不是处处相同，故 D 错误。

故选 B。
}

\item
%% number: 3
%% paperName: 2026年广东高考第 3 题 4 分
%% typeId: 1
%% type: 单选题
%% chapter: 交流电
%% point: 交流电
%% method: 无
%% score: 4
%% degree: 850
%% duplicateId: 0
%% body:
如图是某种自由式活塞内燃发电机输出电压 $U$ 随时间 $t$ 变化的图像。下列说法正确的是 \xzanswer{C}
\onepicture[w=5cm]{figs/03.png}

\fourchoices[answer=C]
{电压的有效值为 $200\sqrt{2}\UV$}
{电压的最大值为 $400\UV$}
{电压的频率为 $50\UHz$}
{电压的周期为 $0.01\Us$}

%% answer: C
%% memo:
\memoanswer{
AB．由图像可知，电压的最大值为 $200\UV$，且并不是一直维持 $200\UV$，则有效值小于 $200\UV$，故 AB 错误；

CD．从 $U-t$ 图像可以看出，电压完成一次周期性变化的时间（周期）$T=0.02\Us$，则频率为 $f=\frac{1}{T}=50\UHz$，故 C 正确，D 错误。

故选 C。
}

\item
%% number: 4
%% paperName: 2026年广东高考第 4 题 4 分
%% typeId: 1
%% type: 单选题
%% chapter: 直线运动
%% point: 多段匀变速直线运动
%% method: 无
%% score: 4
%% degree: 550
%% duplicateId: 0
%% body:
足球比赛中，某队员为接应传球，由静止开始沿直线跑动，先匀加速冲刺，后匀减速至接球点停止。全程用时 $5\Us$，位移大小为 $20\Um$，则该队员在某时刻的速度和加速度的大小可能是 \xzanswer{A}

\fourchoices[answer=A]
{$6\Ums$，$2\Umsq$}
{$6\Ums$，$1\Umsq$}
{$9\Ums$，$2\Umsq$}
{$9\Ums$，$1\Umsq$}

%% answer: A
%% memo:
\memoanswer{
根据题意，设运动过程中最大速度为 $v_{m}$，则有 $x=\frac{v_{m}}{2}t$，代入数据解得 $v_{m}=8\Ums$。根据题意可知，加速过程或减速过程的最大时间均小于总时间，则有 $v_{m}=at$。

解得加速过程或减速过程的最小加速度 $a=\frac{v_{m}}{t}=1.6\Umsq$。

综上所述，该队员某时刻的速度不可能为 $9\Ums$，加速度大小不可能为 $1\Umsq$。所以根据题意和选项，速度和加速度大小可能是 $6\Ums$，$2\Umsq$。

故选 A。
}

\item
%% number: 5
%% paperName: 2026年广东高考第 5 题 4 分
%% typeId: 1
%% type: 单选题
%% chapter: 曲线运动
%% point: 平抛运动
%% method: 无
%% score: 4
%% degree: 550
%% duplicateId: 0
%% body:
如图是月球上一圆柱形阴影坑竖直截面图。假定某飞行器在月面上空 $H=1.28\Ukm$ 向坑中心方向以速度 $v_{0}$ 匀速水平飞行。在距坑边 $l=1.2\Ukm$ 的 $P$ 点正上方关闭动力，此后只受月球重力，直至抵达着陆线。已知坑直径 $d=3\Ukm$，月面至着陆线深度 $h=1.6\Ukm$，月面重力加速度取 $g_{\text{月}}=1.6\Umsq$，飞行器可视为质点。飞行器安全到达着陆线，则 $v_{0}$ 的大小可能是 \xzanswer{B}
\onepicture[w=6cm]{figs/05.png}

\fourchoices[answer=B]
{$20\Ums$}
{$50\Ums$}
{$110\Ums$}
{$120\Ums$}

%% answer: B
%% memo:
\memoanswer{
根据题意可知飞行器安全到达着陆线，速度最小时会刚好通过月面坑边到达着陆线，

此时根据平抛运动规律在竖直方向有 $H=\frac{1}{2}g_{\text{月}}t_{1}^{2}$，代入数据解得 $t_{1}=40\Us$；

在水平方向有 $l=v_{0}t_{1}$，代入数据可得 $v_{0}=30\Ums$；

飞行器安全到达着陆线，速度最大时会到达着陆线最右端的坑边缘，此时根据平抛运动规律在竖直方向有 $H+h=\frac{1}{2}g_{\text{月}}t_{2}^{2}$，代入数据解得 $t_{2}=60\Us$；

在水平方向有 $l+d=v_{0}t_{2}$，代入数据可得 $v_{0}=70\Ums$。

所以飞行器安全到达着陆线需满足 $30\Ums\leq v_{0}\leq70\Ums$。

故选 B。
}

\item
%% number: 6
%% paperName: 2026年广东高考第 6 题 4 分
%% typeId: 1
%% type: 单选题
%% chapter: 曲线运动
%% point: 开普勒定律
%% method: 无
%% score: 4
%% degree: 550
%% duplicateId: 0
%% body:
如图所示，某行星对单个卫星表面最远点与最近点的单位质量物体的“引力差值”可近似为 $F=k\frac{MR}{r^{3}}$，其中 $k$ 为常量，$M$ 为行星质量，$R$ 为卫星球体半径，$r$ 为行星中心到卫星中心的距离。两卫星 P 和 Q 的球体半径之比 $R_{P}:R_{Q}$ 为 $2:1$，它们绕该行星做匀速圆周运动的周期之比 $T_{P}:T_{Q}$ 为 $8:1$，该行星对卫星 P、Q 的“引力差值”分别为 $F_{P}$、$F_{Q}$，则 $F_{P}:F_{Q}$ 为 \xzanswer{C}
\onepicture[w=5cm]{figs/06.png}

\fourchoices[answer=C]
{$1:4$}
{$1:16$}
{$1:32$}
{$1:64$}

%% answer: C
%% memo:
\memoanswer{
根据 $G\frac{Mm}{r^{2}}=m\frac{4\pi^{2}}{T^{2}}r$ 可得 $r^{3}=\frac{GMT^{2}}{4\pi^{2}}$。

又因为两卫星绕该行星做匀速圆周运动的周期之比 $T_{P}:T_{Q}$ 为 $8:1$，可得

$\frac{r_{P}^{3}}{r_{Q}^{3}}=\frac{T_{P}^{2}}{T_{Q}^{2}}=\left(\frac{8}{1}\right)^{2}=64$。

根据“引力差值”公式 $F=k\frac{MR}{r^{3}}$，其中 $M$ 和 $k$ 相同，所以 $\frac{F_{P}}{F_{Q}}=\frac{R_{P}}{R_{Q}}\cdot\frac{r_{Q}^{3}}{r_{P}^{3}}=\frac{2}{1}\times\frac{1}{64}=\frac{1}{32}$。

故选 C。
}

\item
%% number: 7
%% paperName: 2026年广东高考第 7 题 4 分
%% typeId: 1
%% type: 单选题
%% chapter: 电场
%% point: 电容
%% method: 无
%% score: 4
%% degree: 450
%% duplicateId: 0
%% body:
图~\ref{2026广东07a} 为某梳齿状可变电容器截面图，上端为固定极板，下端为可上下运动的动极板。将该电容器接入图~\ref{2026广东07b} 所示电路，探究其电荷量 $Q$、极板电压 $U$ 和电容 $C$ 的变化。过程为：①当动极板运动到最高处，开关 S 接 a 端，电源 $E$ 向电容器充电；②充电结束后 S 接 b 端空置，动极板向下运动；③当动极板运动到最低处，S 接 c 端，电容器通过 $R$ 放电。关于该过程，下列图像可能正确的是 \xzanswer{A}
\twopicture[labelA=2026广东07a, labelB=2026广东07b, h=3.2cm]
{figs/07a.png}
{figs/07b.png}

\fourchoices[answer=A, ispicture=true, h=2.6cm]
{figs/07c.png}
{figs/07d.png}
{figs/07e.png}
{figs/07f.png}

%% answer: A
%% memo:
\memoanswer{
AB．由题意可知①过程为电容器充电过程，电容器两端电势差逐渐增大，根据 $Q=CU$ 可知电容器电荷量逐渐增大；②过程 S 接 b 端，则电容器 $Q$ 不变，动极板向下运动，$d$ 增大，根据 $C=\frac{\varepsilon S}{4\pi kd}=\frac{Q}{U}$ 可知 $U$ 增大，$C$ 减小；③过程为放电过程，动极板在最低处电容为更小的恒定值，故对应 $Q-U$ 图线为过原点、斜率更小的直线，A 正确，B 错误；

CD．通过前面的分析可知①过程的电容 $C$ 大于③过程的电容 $C$，②过程中 $Q$ 不变，因为 $C=\frac{Q}{U}$，所以 $U$ 与 $C$ 成反比，图像为双曲线的一支，故 CD 均错误。

故选 A。
}

\item
%% number: 8
%% paperName: 2026年广东高考第 8 题 6 分
%% typeId: 2
%% type: 多选题
%% chapter: 曲线运动
%% point: 圆周运动的应用
%% method: 无
%% score: 6
%% degree: 650
%% duplicateId: 0
%% body:
如图所示，在光滑的水平地面上，P、Q、M、N 四个质量相等的小球通过两根不可伸长的轻绳相连，P、Q 间的绳长为 $L_{1}$，M、N 间的绳长为 $L_{2}$，两绳相交于各自的中点 O，四球以相同角速度 $\omega$ 绕固定的 O 点做匀速圆周运动，已知 $L_{1}=2\Um$，P、Q 的向心加速度大小均为 $a_{1}=4\Umsq$，M、N 的向心加速度大小均为 $a_{2}=1\Umsq$，四球均可视为质点，忽略空气阻力，下列说法正确的有 \xzanswer{AB}
\onepicture[w=5.5cm]{figs/08.png}

\fourchoices[answer=AB]
{$\omega=2\Urads$}
{$L_{2}=0.5\Um$}
{P 的线速度大小为 $0.5\Ums$}
{轻绳对四球的拉力大小相等}

%% answer: AB
%% memo:
\memoanswer{
A．根据题意可知 P、Q 做圆周运动的半径 $r_{1}=\frac{L_{1}}{2}=1\Um$。

对 P、Q 根据 $a_{1}=\omega^{2}r_{1}$ 可得 $\omega=\sqrt{\frac{a_{1}}{r_{1}}}=\sqrt{\frac{4}{1}}\Urads=2\Urads$，故 A 正确；

B．因为四个球角速度相同，对球 M、N 有 $r_{2}=\frac{L_{2}}{2}=\frac{a_{2}}{\omega^{2}}=\frac{1}{4}\Um$，

所以 $L_{2}=0.5\Um$，故 B 正确；

C．对于 P 有 $v_{1}=\omega r_{1}=1\times2\Ums=2\Ums$，故 C 错误；

D．绳子拉力提供向心力，根据 $F=m\omega^{2}r$ 可知由于四球质量相等，角速度相等，但半径不同，故拉力大小不相等，故 D 错误。

故选 AB。
}

\item
%% number: 9
%% paperName: 2026年广东高考第 9 题 6 分
%% typeId: 2
%% type: 多选题
%% chapter: 光学
%% point: 全反射
%% method: 无
%% score: 6
%% degree: 550
%% duplicateId: 0
%% body:
图中底端为锥形且顶角为 $90^{\circ}$ 的直光纤，可用于检测流动液体中的空气泡，某单色光从光纤顶端左部入射，平行于轴线方向传播，探测器在光纤顶端右部探测锥面反射光的光强，已知锥面外为空气时，该单色光在锥面恰好发生全反射，空气折射率取 $1$，被测液体折射率大于此光纤折射率。下列说法正确的有 \xzanswer{AD}
\onepicture[w=4.3cm]{figs/09.png}

\fourchoices[answer=AD]
{此光纤对该单色光的折射率为 $\sqrt{2}$}
{该单色光在此光纤中的传播速度是真空中光速的 $\frac{1}{2}$}
{锥面浸入液体过程中，探测到的光强相对于浸入液体前变强}
{锥面完全浸入液体后，若探测到的光强在变强，说明检测到气泡}

%% answer: AD
%% memo:
\memoanswer{
A．画出光路图如图所示。
\onepicture[w=4.3cm, align=right]{figs/09_an.png}
根据几何关系可知，该单色光在锥面发生全反射的临界角 $C=45^{\circ}$，

根据临界角公式 $\sin C=\frac{1}{n}$，

可得此光纤对该单色光的折射率为 $n=\sqrt{2}$，故 A 正确；

B．根据 $v=\frac{c}{n}$，可得该单色光在此光纤中的传播速度为 $v=\frac{\sqrt{2}}{2}c$，故 B 错误；

C．光从折射率大的介质到折射率小的介质中，才可能发生全反射，被测液体折射率大于此光纤折射率，则锥面浸入液体部分，光不会发生全反射，即锥面浸入液体过程中，探测到的光强相对于浸入液体前变弱，故 C 错误；

D．锥面完全浸入液体后，若探测到的光强在变强，说明光在某些地方发生了全反射，发生全反射的地方即有空气，说明检测到气泡，故 D 正确。

故选 AD。
}

\item
%% number: 10
%% paperName: 2026年广东高考第 10 题 6 分
%% typeId: 2
%% type: 多选题
%% chapter: 磁场
%% point: 质谱仪回旋加速器
%% method: 无
%% score: 6
%% degree: 350
%% duplicateId: 0
%% body:
如图是一种长方体电子磁谱仪结构示意图，磁谱仪内存在磁感应强度大小为 $B$、方向垂直底面向上的匀强磁场，磁场区域长为 $a$、宽为 $b$。电子束中有三个电子通过准直器后垂直左侧面沿边缘进入磁场，偏转后分别到达磁谱仪三个侧面，与边缘的距离分别为 $x_{1}$、$x_{2}$ 和 $x_{3}$，电子电荷量为 $-e$、质量为 $m$，不考虑相对论效应，下列说法正确的有 \xzanswer{BCD}
\onepicture[w=7cm]{figs/10.png}

\fourchoices[answer=BCD]
{电子 1 的动能比电子 3 的大}
{电子 1 在磁场中的运动时间为 $\frac{\pi m}{eB}$}
{电子 2 的动能为 $\frac{e^{2}B^{2}(x_{2}^{2}+b^{2})^{2}}{8mb^{2}}$}
{电子 3 的动量大小为 $\frac{eB(x_{3}^{2}+a^{2})}{2x_{3}}$}

%% answer: BCD
%% memo:
\memoanswer{
A．电子在磁场中做匀速圆周运动，根据洛伦兹力提供向心力可得 $evB=m\frac{v^{2}}{r}$，

可得 $v=\frac{eBr}{m}$。

根据几何关系可知电子 1 的运动半径小于电子 3 的运动半径，则电子 1 的速度小于电子 3 的速度，根据 $E_{k}=\frac{1}{2}mv^{2}$ 可知电子 1 的动能比电子 3 的小，故 A 错误；

B．电子在磁场中做匀速圆周运动的周期 $T=\frac{2\pi r}{v}$，结合 $v=\frac{eBr}{m}$，可得 $T=\frac{2\pi m}{eB}$。

根据题图可知电子 1 在磁场中运动的圆心角为 $180^{\circ}$，则电子 1 在磁场中的运动时间为 $t=\frac{180^{\circ}}{360^{\circ}}\times T=\frac{\pi m}{eB}$，故 B 正确；

C．设电子 2 运动半径为 $R_{2}$，如图所示，
\onepicture[w=5cm]{figs/10_an1.png}
根据几何关系有 $(b-R_{2})^{2}+x_{2}^{2}=R_{2}^{2}$，

可得 $R_{2}=\frac{b^{2}+x_{2}^{2}}{2b}$。

根据 $v_{2}=\frac{eBR_{2}}{m}$ 可得 $v_{2}=\frac{eB(x_{2}^{2}+b^{2})}{2bm}$，

电子 2 的动能为 $E_{k}=\frac{1}{2}mv_{2}^{2}=\frac{e^{2}B^{2}(x_{2}^{2}+b^{2})^{2}}{8mb^{2}}$，故 C 正确；

D．设电子 3 运动半径为 $R_{3}$，如图所示，
\onepicture[w=5cm]{figs/10_an2.png}
根据几何关系有 $(R_{3}-x_{3})^{2}+a^{2}=R_{3}^{2}$，

可得 $R_{3}=\frac{x_{3}^{2}+a^{2}}{2x_{3}}$。

根据 $v_{3}=\frac{eBR_{3}}{m}$ 可得 $v_{3}=\frac{eB(x_{3}^{2}+a^{2})}{2mx_{3}}$，

电子 3 的动量大小为 $p_{3}=mv_{3}=\frac{eB(x_{3}^{2}+a^{2})}{2x_{3}}$，故 D 正确。

故选 BCD。
}

\item
%% number: 11
%% paperName: 2026年广东高考第 11 题 8 分
%% typeId: 4
%% type: 实验题
%% chapter: 直线运动
%% point: 打点计时器公式
%% method: 无
%% score: 8
%% degree: 550
%% duplicateId: 0
%% body:
某科技小组设计了测量薄膜压缩时间的实验，图~\ref{2026广东11a} 为装置示意图。
\twopicture[labelA=2026广东11a, labelB=2026广东11b, h=4.5cm]
{figs/11a.png}
{figs/11b.png}

（1）实验原理。

某种薄膜受到小球冲击时会发光，通过测量发光光谱的相对光强峰值 $I_{p}$，可得最大冲击力 $F$，结合动量定理可估测薄膜压缩时间 $t$。

（2）实验操作及数据处理。

\begin{steps}
	\item
	称量小球质量，记录质量 $m$。
	\item
	断开电磁铁电源，使小球从薄膜正上方自由下落。拍摄小球下落过程的视频。
	\item
	图~\ref{2026广东11b} 为利用视频处理软件得到的小球接触薄膜瞬间前连续 $3$ 个时刻的位置图。分别测量图中的距离 $s_{1}$ 和 $s_{2}$，可得 $s_{1}$ 为 \tkanswer{6.26} $\Ucm$，$s_{2}$ 为 $6.66\Ucm$。相邻两个位置的时间间隔为 $T$，则小球在 B 位置处的速度 $v_{1}=$ \tkanswer{$\frac{s_{1}+s_{2}}{2T}$}（用 $s_{1}$、$s_{2}$ 和 $T$ 表示）；与薄膜接触前瞬间，小球在 C 位置处的速度 $v_{2}=$ \tkanswer{$\frac{3s_{2}-s_{1}}{2T}$}（用 $s_{1}$、$s_{2}$ 和 $T$ 表示）。
	\item
	用光谱仪测得薄膜受到冲击时发光光谱的 $I_{p}$ 为 $158$ a.u.。图~\ref{2026广东11c} 为已知的 $F-I_{p}$ 关系图像，可读得 $F$ 为 \tkanswer{100} $\UN$。
	\item
	$m=30.0\Ug$，$T=0.02\Us$，由动量定理可估测薄膜压缩时间 $t=\frac{2mv_{2}}{F}$，将数据代入可得 $t=$ \tkanswer{$2.06\times10^{-3}$} $\Us$（结果保留 3 位有效数字）。
\end{steps}
\onepicture[num=c, label=2026广东11c, showlabel=false, align=right, w=5.5cm]{figs/11c.png}

%% answer: ③ 6.26，$\frac{s_{1}+s_{2}}{2T}$，$\frac{3s_{2}-s_{1}}{2T}$；④ $100\UN$；⑤ $2.06\times10^{-3}\Us$
%% memo:
\memoanswer{
①毫米刻度尺的分度值为 $1\Umm$，需要估读到下一位，根据题图可得 $s_{1}=262.6\Umm-200.0\Umm=62.6\Umm=6.26\Ucm$；

②根据匀变速直线运动中间时刻瞬时速度等于该过程平均速度可得 $v_{1}=\frac{s_{1}+s_{2}}{2T}$；

③根据速度时间关系有 $v_{2}=v_{1}+gT$，又 $s_{2}-s_{1}=gT^{2}$，联立可得 $v_{2}=\frac{3s_{2}-s_{1}}{2T}$；

④根据题图可得 $I_{p}$ 为 $158$ a.u. 时，$F$ 为 $100\UN$；

⑤$m=30.0\Ug$，$T=0.02\Us$，由动量定理可估测薄膜压缩时间 $t=\frac{2mv_{2}}{F}$，结合 $v_{2}=\frac{3s_{2}-s_{1}}{2T}$，将数据代入可得 $t=2.06\times10^{-3}\Us$。
}

\item
%% number: 12
%% paperName: 2026年广东高考第 12 题 8 分
%% typeId: 4
%% type: 实验题
%% chapter: 电路
%% point: 传感器
%% method: 无
%% score: 8
%% degree: 550
%% duplicateId: 0
%% body:
棉花的回潮率可通过测量一定压力下棉花的电阻得到。某科技小组制作了利用该方法测量棉花回潮率的简易装置，如图~\ref{2026广东12a}。旋转螺杆压缩弹簧对棉花施加压力。由图~\ref{2026广东12b} 所示电路测量棉花的电阻 $R_{c}$，所用器材有：电源 $E$；微安表 $\mu$A；定值电阻 $R_{0}$；电阻箱 $R$；开关 $S_{1}$ 和 $S_{2}$；导线若干。
\threepicture[labelA=2026广东12a, labelB=2026广东12b, labelC=2026广东12c, h=3.4cm]
{figs/12a.png}
{figs/12b.png}
{figs/12c.png}

（1）弹簧劲度系数测量。

\begin{steps}
	\item
	将装有弹簧的绝缘压板放到水平桌面上，如图~\ref{2026广东12c}。用刻度尺测量并记录弹簧原长。
	\item
	在弹簧上端加一个砝码，待砝码 \tkanswer[5cm]{平衡稳定（或静止，表述合理即可）} 后读数，记录弹簧长度。
	\item
	依次增加砝码，重复步骤②。
	\item
	根据实验数据，描绘压力 $F$ 与弹簧压缩量 $x$ 的 $F-x$ 图线，图线为直线，则图线的 \tkanswer{斜率} 表示弹簧劲度系数。
\end{steps}

（2）棉花的电阻 $R_{c}$ 测量。

\begin{steps}
	\item
	称量装有弹簧和刻度尺的绝缘压板质量，记录质量。
	\item
	将某棉花样品装进装置中，装上绝缘压板，确保环形电极与棉花接触良好，固定顶盖，正确连接电路。
	\item
	闭合 $S_{1}$ 和 $S_{2}$，调节 $R$，使 $\mu$A 指针偏转到某一合适位置。此时 $R$ 的示数如图~\ref{2026广东12d}，读数为 \tkanswer{$29\,000$} $\UO$，记为 $R_{1}$。断开 $S_{1}$ 和 $S_{2}$。
	\item
	缓慢旋进螺杆压紧棉花，确保压力达到回潮率测量的要求。$S_{2}$ 保持断开，闭合 $S_{1}$，发现 $\mu$A 示数变 \tkanswer{小}，将 $R$ 的阻值 \tkanswer{调小}，使 $\mu$A 指针偏转到与步骤③相同的位置。记录 $R$ 的示数 $R_{2}$。断开 $S_{1}$，利用所测量的物理量，可得棉花的电阻表达式 $R_{c}=$ \tkanswer{$R_{1}-R_{2}$}。
\end{steps}

（3）利用定标曲线确定棉花的回潮率。
\onepicture[num=d, label=2026广东12d, showlabel=false, align=right, w=4cm]{\includesvg{figs/12d}}

%% answer: （1）② 平衡稳定（或静止，表述合理即可）；④ 斜率。（2）③ $29\,000\UO$；④ 小，调小，$R_{1}-R_{2}$
%% memo:
\memoanswer{
（1）添加砝码后，需要等待砝码平衡稳定后再读数，才能保证弹簧弹力等于砝码重力，得到准确的弹簧长度。

根据胡克定律 $F=kx$，因此 $F-x$ 图线的斜率等于劲度系数 $k$。

（2）电阻箱读数为各转盘读数乘倍率之和：$R=2\times10000+9\times1000\UO=29000.0\UO$。

$S_{2}$ 断开，$R_{c}$ 串联接入电路，总电阻增大，因此电流减小，微安表示数变小。

要让电流回到原来的数值，需要总电阻和原来相等，因此需要调小电阻箱 $R$ 的阻值，使总电阻回到原来大小。

等效后满足 $R_{1}=R_{2}+R_{c}$，整理得棉花电阻 $R_{c}=R_{1}-R_{2}$。
}

\item
%% number: 13
%% paperName: 2026年广东高考第 13 题 9 分
%% typeId: 6
%% type: 计算题
%% chapter: 气体性质
%% point: 玻意耳定律
%% method: 无
%% score: 9
%% degree: 550
%% duplicateId: 0
%% body:
图~\ref{2026广东13a} 所示的空气垫是由多个相连的独立气室构成的包装材料，其简化模型如图~\ref{2026广东13b}。充气前气室内均没有气体，在室温 $T_{0}$ 下，将压强 $p_{0}$、体积 $V_{0}$ 的气体通过单向阀充入 $10$ 个气室（忽略气道内气体），此时每个气室均为圆柱体，横截面半径为 $r$，长度为 $h$，当充气后的气室受到挤压变形时，其横截面变成图~\ref{2026广东13c} 所示的“跑道”形（两端是直径为 $d$ 的半圆），且气室长度、横截面周长均保持不变，气室内气体可视为理想气体，充气及挤压变形过程中气体温度始终与室温相同。
\begin{enumerate}
	\item
	求充气后未挤压变形时气室中的压强 $p_{1}$；
	\item
	求挤压变形后气室中的压强 $p_{2}$；
	\item
	已知气室中的压强超过 $p_{c}$ 时气室会爆破，若气室经如图~\ref{2026广东13c} 所示的挤压变形后，体积不变、室温升高，求气室不爆破的最高室温 $T_{c}$。
\end{enumerate}
\threepicture[labelA=2026广东13a, labelB=2026广东13b, labelC=2026广东13c, h=2.6cm, align=right]
{figs/13a.png}
{figs/13b.png}
{figs/13c.png}

%% answer:
\jdanswer{
\begin{enumerate}
	\item $p_{1}=\frac{p_{0}V_{0}}{10\pi hr^{2}}$
	\item $p_{2}=\frac{2p_{0}V_{0}}{5\pi dh(4r-d)}$
	\item $T_{c}=\frac{5\pi dh(4r-d)p_{c}T_{0}}{2p_{0}V_{0}}$
\end{enumerate}
}
%% memo:
\memoanswer{
（1）充气过程温度不变，对所有充入气室的气体，由玻意耳定律有 $p_{0}V_{0}=p_{1}V_{1}$。

根据题意有 $V_{1}=10\pi r^{2}\cdot h$。

解得充气后未挤压变形时气室中的压强 $p_{1}=\frac{p_{0}V_{0}}{10\pi hr^{2}}$。

（2）变形前，气室的体积为 $V=\pi r^{2}h$。

挤压过程温度不变，且横截面周长不变，原圆形横截面周长 $C=2\pi r$。

变形后跑道形横截面两端半圆合为一个整圆，周长为 $C_{1}=\pi d$。

剩余直边总长度为 $C_{2}=C-C_{1}=2\pi r-\pi d$。

单根直边长 $x=\frac{C_{2}}{2}=\pi r-\frac{\pi d}{2}$。

变形后，气室的体积为 $V_{2}=\left[\pi\left(\frac{d}{2}\right)^{2}+xd\right]\cdot h$。

对气室气体，等温过程由玻意耳定律 $p_{1}V=p_{2}V_{2}$。

解得挤压变形后气室中的压强 $p_{2}=\frac{2p_{0}V_{0}}{5\pi dh(4r-d)}$。

（3）体积不变，气体等容变化，有 $\frac{p_{2}}{T_{0}}=\frac{p_{c}}{T_{c}}$。

解得气室不爆破的最高室温 $T_{c}=\frac{5\pi dh(4r-d)p_{c}T_{0}}{2p_{0}V_{0}}$。
}

\item
%% number: 14
%% paperName: 2026年广东高考第 14 题 13 分
%% typeId: 6
%% type: 计算题
%% chapter: 运动和力
%% point: 动量和能量
%% method: 无
%% score: 13
%% degree: 450
%% duplicateId: 0
%% body:
如图是一种球形机器人跳跃原理的示意图，水平横轴过球心 $O$ 点与外壳固定，外壳上的两挡板位于过 $O$ 点的水平线上，两质量均为 $m$ 的摆锤，由长均为 $R$ 的不可伸长轻绳悬挂于轴上的 $O$ 点，初始时刻，两摆锤同时以水平初速度 $v_{0}$ 从最低点向相反方向摆动，直至与两挡板发生碰撞，碰撞时间极短，随后带动外壳以共同速度竖直向上运动，机器人到达最高点后落回地面瞬间，外壳立即静止，两摆锤速度不变，与挡板分离，继续向下运动，已知机器人（含摆锤）总质量为 $M=4\Ukg$，$m=1\Ukg$，$R=0.4\Um$，$v_{0}=4\Ums$。重力加速度取 $g=10\Umsq$，忽略空气阻力，摆锤可视为质点，求：
\begin{enumerate}
	\item
	摆锤与挡板碰撞后瞬间，机器人的动能 $E_{k}$；
	\item
	机器人外壳上升的最大高度 $h$；
	\item
	从摆锤开始运动到第一次外壳落地静止过程中的机械能损失 $\Delta E$。
\end{enumerate}
\onepicture[w=3.8cm, align=right]{figs/14.png}

%% answer:
\jdanswer{
\begin{enumerate}
	\item $E_{k}=4\UJ$
	\item $h=0.1\Um$
	\item $\Delta E=6\UJ$
\end{enumerate}
}
%% memo:
\memoanswer{
（1）两摆锤以初速度 $v_{0}=4\Ums$ 沿外壳向上运动，与挡板相碰前，摆锤机械能守恒。

设摆锤与挡板相碰前的速度为 $v$，根据机械能守恒定律有 $\frac{1}{2}mv_{0}^{2}=mgR+\frac{1}{2}mv^{2}$，

解得 $v=2\sqrt{2}\Ums$。

摆锤与挡板相碰后与机器人一起运动，根据动量守恒定律有 $2mv=Mv_{\text{共}}$，

解得 $v_{\text{共}}=\sqrt{2}\Ums$。

机器人起跳时的动能 $E_{k}=\frac{1}{2}Mv_{\text{共}}^{2}=4\UJ$。

（2）根据速度位移关系 $0-v_{\text{共}}^{2}=-2gh$，

可得机器人外壳上升的最大高度 $h=0.1\Um$。

（3）机器人外壳落到地面时，机器人外壳的速度立即变为 $0$，根据竖直上抛运动的对称性可知摆锤速度大小为 $\sqrt{2}\Ums$，从摆锤开始运动到第一次外壳落地静止过程中的机械能损失 $\Delta E=2\times\frac{1}{2}mv_{0}^{2}-\left(2\times\frac{1}{2}mv_{\text{共}}^{2}+2mgR\right)$，

解得 $\Delta E=6\UJ$。
}

\item
%% number: 15
%% paperName: 2026年广东高考第 15 题 16 分
%% typeId: 6
%% type: 计算题
%% chapter: 电场
%% point: 带电粒子的往复运动
%% method: 无
%% score: 16
%% degree: 350
%% duplicateId: 0
%% body:
图~\ref{2026广东15a} 所示，两竖直放置且足够大的平行金属板 M、N，两板间距为 $d$，在两板正中间竖直平面内固定有一水平绝缘横杆，一质量为 $m$、电荷量为 $q$ 的小球通过两根等长且不可伸长的绝缘轻绳悬挂于横杆下方，小球与横杆的距离为 $d$，两绳的夹角为直角，如图~\ref{2026广东15b}，接通电源，使板间电压由 $0$ 开始缓慢增大，小球缓慢向 N 靠近，在此过程中每个时刻小球都受力平衡，当小球接触 N 的瞬间，电荷量变为 $-q$，板间电压停止增大并在此后保持恒定，在此恒定电压下，小球每次与 M 或 N 接触后瞬间，速度均减为 $0$，带电荷量变化满足“电性反转、大小不变”，从而在两板间沿着圆弧往复运动，重力加速度为 $g$，小球可视为质点，每次与板碰撞均不影响两板间电压，忽略空气阻力和电场的边缘效应，忽略小球所带电荷对板间电场的影响。
\begin{enumerate}
	\item
	求 M、N 间的恒定电压 $U$；
	\item
	求小球第一次碰撞 M 前瞬间，单根轻绳的拉力大小 $T$；
	\item
	若某次小球碰撞 M 时，M、N 间的电压突变为原恒定电压的 $k$ 倍（$k>0$），其他条件不变，此后小球仍能沿着圆弧往复运动，求 $k$ 的取值范围，并求出该范围内不同 $k$ 值对应的小球最大动能 $E_{k}$。
\end{enumerate}
\twopicture[labelA=2026广东15a, labelB=2026广东15b, h=4cm, align=right]
{figs/15a.png}
{figs/15b.png}

%% answer:
\jdanswer{
\begin{enumerate}
	\item $U=\frac{\sqrt{3}mgd}{3q}$
	\item $T=\frac{2\sqrt{6}}{3}mg$
	\item 当 $0<k\leq1$ 时，$E_{k}=\left(\frac{\sqrt{3}}{6}k+\sqrt{1+\frac{k^{2}}{3}}-\frac{\sqrt{3}}{2}\right)mgd$；当 $1<k<3$ 时，$E_{k}=\frac{\sqrt{3}}{3}kmgd$。
\end{enumerate}
}
%% memo:
\memoanswer{
（1）当小球刚好到金属板 N 时，受力分析如图所示。
\onepicture[w=4.2cm]{figs/15_an1.png}
根据几何关系可得 $\sin\theta=\frac{\frac{d}{2}}{d}$，可得 $\theta=30^{\circ}$。

根据平衡条件有 $qE=mg\tan\theta$，其中 $E=\frac{U}{d}$。

联立可得 $U=\frac{\sqrt{3}mgd}{3q}$。

（2）小球从金属板 N 到金属板 M 过程，根据动能定理
\onepicture[w=4.2cm]{figs/15_an2.png}
$qU=\frac{1}{2}mv^{2}-0$。

小球第一次碰撞 M 前瞬间，受力分析如图所示。

$2T\cos45^{\circ}-mg\cos\theta-qE\sin\theta=m\frac{v^{2}}{d}$。

对小球根据牛顿第二定律，

联立解得 $T=\frac{2\sqrt{6}}{3}mg$。

（3）如图所示，要使小球仍能沿着圆弧往复运动，即绳子拉力大于零，需要满足 $q\cdot\frac{kU}{d}\sin\theta<mg\cos\theta$，
\onepicture[w=4.2cm]{figs/15_an3.png}
解得 $k\leq3$。

根据（2）分析可知当 $k=1$ 时，小球到达金属板 N 时，电场力和重力的合力与绳子的合力方向相同，当 $1<k<3$ 时，小球所受的重力和电场力的合力相对于小球到达金属板 N 时绳子合力的方向偏上，如图所示。
\onepicture[w=4.2cm]{figs/15_an4.png}
绳子拉力对小球不做功，可知小球所受的重力和电场力的合力与小球位移的夹角一直为锐角，当小球到达金属板 N 时，动能最大，根据动能定理有 $kqU=E_{k}-0$，

可得 $E_{k}=\frac{\sqrt{3}}{3}kmgd$。

当 $0<k\leq1$ 时，小球所受的重力和电场力的合力相对于小球到达金属板 N 时绳子合力的方向偏下，如图所示。
\onepicture[w=4.2cm]{figs/15_an5.png}
当小球速度与小球所受的重力和电场力的合力方向垂直时，速度最大，该点为等效最低点，设此时与竖直方向的夹角为 $\alpha$，根据几何关系有

$\tan\alpha=\frac{\frac{kU}{d}\cdot q}{mg}=\frac{\sqrt{3}}{3}k$。

根据数学知识可得

$\sin\alpha=\frac{\frac{\sqrt{3}}{3}k}{\sqrt{1+\frac{1}{3}k^{2}}}$，$\cos\alpha=\frac{1}{\sqrt{1+\frac{1}{3}k^{2}}}$。

根据动能定理有

$mgd(\cos\alpha-\cos\theta)+\frac{kU}{d}\cdot qd(\sin\theta+\sin\alpha)=E_{k}-0$。

解得 $E_{k}=\left(\frac{\sqrt{3}}{6}k+\sqrt{1+\frac{k^{2}}{3}}-\frac{\sqrt{3}}{2}\right)mgd$。
}

\end{enumerate}

\end{document}
